% !TEX program = xelatex
% Lernplatt - by Can Siebert
% Kurs 22 (Bo 143) - Tag 10 - Loesungsheft
% Performance-Marketing im Vertrieb: Anzeigen, Conversion & E-Mail-Kampagnen
%
% Self-contained xelatex source. Kein biblatex, keine externen Bilder.

\documentclass[11pt,a4paper,oneside]{article}

% --- Encoding & Sprache --------------------------------------------------
\usepackage{polyglossia}
\setdefaultlanguage{german}

% --- Geometrie -----------------------------------------------------------
\usepackage[a4paper,margin=2.5cm]{geometry}

% --- Fonts ---------------------------------------------------------------
\usepackage{fontspec}
\defaultfontfeatures{Ligatures=TeX}

\IfFontExistsTF{Charter}{%
  \setmainfont{Charter}%
}{%
  \IfFontExistsTF{Noto Serif}{%
    \setmainfont{Noto Serif}%
  }{%
    \setmainfont{Liberation Serif}%
  }%
}

\IfFontExistsTF{Source Sans 3}{%
  \setsansfont{Source Sans 3}%
}{%
  \IfFontExistsTF{Source Sans Pro}{%
    \setsansfont{Source Sans Pro}%
  }{%
    \setsansfont{Liberation Sans}%
  }%
}

\newcommand{\caveatfontfallback}{\itshape}
\IfFontExistsTF{Caveat}{%
  \newfontfamily\caveatfont{Caveat}[Scale=1.15]%
}{%
  \newcommand\caveatfont{\caveatfontfallback}%
}

\setmonofont{Liberation Mono}

% --- Mikro-Typografie ----------------------------------------------------
\usepackage{microtype}
\usepackage{eurosym}
\linespread{1.15}

% --- Farben & Boxen ------------------------------------------------------
\usepackage{xcolor}
\definecolor{lpInk}{RGB}{20,28,38}
\definecolor{lpAccent}{RGB}{22,90,140}
\definecolor{lpAccentLight}{RGB}{210,228,240}
\definecolor{lpAmber}{RGB}{222,138,30}
\definecolor{lpAmberLight}{RGB}{255,243,217}
\definecolor{lpAmberBorder}{RGB}{196,118,18}
\definecolor{lpGray}{RGB}{96,104,114}
\definecolor{lpGrayLight}{RGB}{238,240,243}
\definecolor{lpGreen}{RGB}{34,120,72}
\definecolor{lpGreenLight}{RGB}{222,238,228}
\definecolor{lpGreenBorder}{RGB}{24,96,58}

\definecolor{lpL1}{RGB}{34,120,72}
\definecolor{lpL2}{RGB}{22,90,140}
\definecolor{lpL3}{RGB}{170,42,52}

\usepackage[most]{tcolorbox}
\tcbuselibrary{breakable,skins}

\newtcolorbox{infobox}[1][Hinweis]{%
  enhanced, breakable,
  colback=lpAccentLight,
  colframe=lpAccent,
  boxrule=0.6pt,
  arc=2pt,
  left=10pt,right=10pt,top=8pt,bottom=8pt,
  fonttitle=\sffamily\bfseries\color{white},
  coltitle=white,
  title={#1},
  attach boxed title to top left={xshift=8pt,yshift=-8pt},
  boxed title style={size=small, colback=lpAccent, colframe=lpAccent, boxrule=0.4pt, arc=1pt},
  top=14pt
}

\newtcolorbox{loesungbox}[1][Musterlösung]{%
  enhanced, breakable,
  colback=lpGreenLight,
  colframe=lpGreenBorder,
  boxrule=0.6pt,
  arc=2pt,
  left=10pt,right=10pt,top=8pt,bottom=8pt,
  fonttitle=\sffamily\bfseries\color{white},
  coltitle=white,
  title={#1},
  attach boxed title to top left={xshift=8pt,yshift=-8pt},
  boxed title style={size=small, colback=lpGreenBorder, colframe=lpGreenBorder, boxrule=0.4pt, arc=1pt},
  top=14pt
}

% --- Listen --------------------------------------------------------------
\usepackage{enumitem}
\setlist{itemsep=2pt, topsep=4pt, parsep=0pt}
\setlist[itemize,1]{label={\textbullet},leftmargin=1.6em}
\setlist[itemize,2]{label={\textendash},leftmargin=1.4em}
\setlist[enumerate,1]{label=\arabic*., leftmargin=1.8em}

% --- Tabellen ------------------------------------------------------------
\usepackage{tabularx}
\usepackage{array}
\usepackage{booktabs}
\usepackage{longtable}

% --- Kopf-/Fusszeilen ----------------------------------------------------
\usepackage{fancyhdr}
\usepackage{lastpage}
\pagestyle{fancy}
\fancyhf{}
\renewcommand{\headrulewidth}{0.3pt}
\renewcommand{\footrulewidth}{0pt}
\fancyhead[L]{\footnotesize\sffamily\color{lpGray}Lösungsheft \textperiodcentered{} Kurs 22 \textperiodcentered{} Tag 10}
\fancyhead[R]{\footnotesize\sffamily\color{lpGray}E-Mail \& Performance-Marketing}
\fancyfoot[L]{\footnotesize\sffamily\color{lpGray}\textcopyright{} Can Siebert}
\fancyfoot[R]{\footnotesize\sffamily\color{lpGray}\thepage{} / \pageref{LastPage}}

\fancypagestyle{plain}{%
  \fancyhf{}%
  \renewcommand{\headrulewidth}{0pt}%
  \fancyfoot[C]{\footnotesize\sffamily\color{lpGray}\thepage{} / \pageref{LastPage}}%
}

% --- Section-Styling -----------------------------------------------------
\usepackage[explicit]{titlesec}

\titleformat{\section}[block]
  {\sffamily\Large\bfseries\color{lpAccent}}
  {\thesection.}{0.6em}{#1}
  [{\vspace{2pt}\color{lpAccent}\titlerule[0.6pt]}]

\titleformat{\subsection}[block]
  {\sffamily\large\bfseries\color{lpInk}}
  {\thesubsection}{0.6em}{#1}

\titlespacing*{\section}{0pt}{14pt}{8pt}
\titlespacing*{\subsection}{0pt}{10pt}{4pt}

% --- Hyperlinks ----------------------------------------------------------
\usepackage[hidelinks,unicode]{hyperref}

% --- Helfer --------------------------------------------------------------
\newcommand{\akzent}[1]{{\caveatfont\color{lpAmber}#1}}
\newcommand{\fachbegriff}[1]{\textbf{#1}}

\newcommand{\niveaubadge}[2]{%
  \tcbox[on line, colback=#1, colframe=#1, coltext=white,
         boxrule=0pt, arc=2pt, left=5pt,right=5pt,top=1pt,bottom=1pt,
         nobeforeafter]{\sffamily\bfseries\footnotesize #2}%
}
\newcommand{\nivLi}{\niveaubadge{lpL1}{L1\,\textbullet\,Wissen}}
\newcommand{\nivLii}{\niveaubadge{lpL2}{L2\,\textbullet\,Anwendung}}
\newcommand{\nivLiii}{\niveaubadge{lpL3}{L3\,\textbullet\,Management}}

\newcommand{\zeit}[1]{%
  \tcbox[on line, colback=lpGrayLight, colframe=lpGrayLight, coltext=lpInk,
         boxrule=0pt, arc=2pt, left=5pt,right=5pt,top=1pt,bottom=1pt,
         nobeforeafter]{\sffamily\footnotesize\textbf{$\sim$\,#1}}%
}

\newcommand{\aufgabenkopf}[4]{%
  \par\vspace{6pt}
  \noindent{\sffamily\Large\bfseries\color{lpAccent}Aufgabe #1 \textendash{} #2}\par
  \vspace{2pt}
  \noindent{\color{lpAccent}\rule{\linewidth}{0.6pt}}\par
  \vspace{4pt}
  \noindent #3 \hfill \zeit{#4}\par
  \vspace{6pt}
}

\newcommand{\ytquery}[1]{%
  \par\vspace{4pt}
  \noindent{\footnotesize\sffamily\color{lpGray}\textbf{Lernvideo zum Thema:}\ %
    \href{https://studyflix.de/search?query=#1}{\textcolor{lpAccent}{Studyflix-Treffer}}\ ·\ %
    \href{https://www.youtube.com/results?search\_query=#1}{\textcolor{lpAccent}{YouTube-Treffer}}\\%
    \textit{\textcolor{lpGray}{Stichworte: #1}}}\par
}

\newcommand{\ytvideo}[2]{%
  \par\vspace{4pt}
  \noindent{\footnotesize\sffamily\color{lpGray}\textbf{Lernvideo:}\ \href{#2}{\textcolor{lpAccent}{#1}}}\par
}

% =========================================================================
\begin{document}

% --- COVER ---------------------------------------------------------------
\begin{titlepage}
  \pagestyle{empty}
  \vspace*{1.5cm}
  \noindent{\sffamily\footnotesize\color{lpGray}LERNPLATT \textperiodcentered{} BY CAN SIEBERT}\\[2pt]
  \noindent{\color{lpAccent}\rule{\linewidth}{1.2pt}}\\[4pt]
  \noindent{\sffamily\footnotesize\color{lpGray}LÖSUNGSHEFT \textperiodcentered{} MODUL 5 \textperiodcentered{} TAG 10 \textperiodcentered{} KURS 22 (Bo 143) \textperiodcentered{} 11.\,Juni 2026}

  \vspace{3cm}
  {\sffamily\fontsize{30}{34}\selectfont\bfseries\color{lpInk}Lösungsheft\\Tag 10\par}

  \vspace{0.7cm}
  {\sffamily\large\color{lpGray}Musterlösungen zu den elf Aufgaben\\des Aufgabenhefts \textendash{} Performance-\\Marketing, E-Mail \& Conversion.\par}

  \vspace{1.5cm}
  {\caveatfont\LARGE\color{lpGreen}Musterlösungen \textperiodcentered{} nicht die einzige Antwort}\par

  \vfill

  \noindent\begin{minipage}{0.55\linewidth}
    {\sffamily\footnotesize\color{lpGray}Hinweis:}\\
    {\sffamily\small\color{lpInk}Musterlösungen sind Diskussionsbasis,\\nicht die einzige korrekte Lösung.}\\[10pt]
    {\sffamily\footnotesize\color{lpGray}Begleitend:}\\
    {\sffamily\small\color{lpInk}Aufgabenheft \& Lehrskript Tag 10}
  \end{minipage}\hfill
  \begin{minipage}{0.4\linewidth}
    \raggedleft
    {\sffamily\footnotesize\color{lpGray}Dozent:}\\
    {\sffamily\small\color{lpInk}Can Siebert}\\[10pt]
    {\sffamily\footnotesize\color{lpGray}Niveau-Mix:}\\
    {\sffamily\small\color{lpInk}3\,L1 \textperiodcentered{} 5\,L2 \textperiodcentered{} 3\,L3}
  \end{minipage}

  \vspace{0.5cm}
  \noindent{\color{lpAccent}\rule{\linewidth}{0.6pt}}
\end{titlepage}

% --- Anleitung -----------------------------------------------------------
\section*{Hinweise zu den Musterlösungen}
\thispagestyle{plain}

\noindent Die folgenden Musterlösungen sind \emph{eine} mögliche, didaktisch nachvollziehbare Lösung \textendash{} nicht die einzig richtige. Insbesondere auf L2 und L3 sind mehrere Begründungswege legitim, solange sie:

\begin{itemize}
  \item den im Lehrskript eingeführten Begriffsapparat (Nurturing, MQL/SQL, ROI, CAC) sauber nutzen,
  \item Annahmen explizit machen, statt sie zu verschleiern,
  \item Zielkonflikte benennen, statt sie wegzudefinieren,
  \item zu einer klar formulierten Entscheidung führen \textendash{} nicht bloß zu einer Beschreibung.
\end{itemize}

\begin{infobox}[Nutzung im Coaching]
Vergleichen Sie Ihre eigene Antwort \emph{vor} der Musterlösung. Wo weicht Ihre Argumentation ab? Ist die Abweichung eine andere Annahme \textendash{} oder ein Denkfehler? Genau dort entsteht der Lerneffekt.
\end{infobox}

\clearpage

% =========================================================================
% L1 AUFGABEN
% =========================================================================
\section*{Niveau L1 \textendash{} Wissen \& Reproduktion}
\addcontentsline{toc}{section}{L1 -- Wissen}

% --- AUFGABE 1 -----------------------------------------------------------
% LERNPLATT_HILFSVIDEOS
\section*{Hilfsvideos zu diesem Tag}
\addcontentsline{toc}{section}{Hilfsvideos zu diesem Tag}
\thispagestyle{plain}

\noindent Die folgenden YouTube-Erkl\"arvideos vermitteln die Inhalte, die Sie zur Bearbeitung der Aufgaben ben\"otigen. Klicken Sie auf den Titel, um das Video direkt zu \"offnen; die vollst\"andige URL steht in Klammern.

\begin{description}[style=nextline,leftmargin=18pt,labelindent=0pt,itemsep=8pt]
  \item[1.~E-Mail-Marketing als systematischer Vertriebshebel] \href{https://youtu.be/GlKFLDx0Fjw}{\textcolor{lpAccent}{https://youtu.be/GlKFLDx0Fjw}}\\[2pt]
  {\footnotesize\color{lpGray}(URL: \nolinkurl{https://youtu.be/GlKFLDx0Fjw})}
  \item[2.~Lead Nurturing mit Marketing Automation] \href{https://youtu.be/sjAkl8pwLgg}{\textcolor{lpAccent}{https://youtu.be/sjAkl8pwLgg}}\\[2pt]
  {\footnotesize\color{lpGray}(URL: \nolinkurl{https://youtu.be/sjAkl8pwLgg})}
  \item[3.~Newsletter erfolgreich starten — Aufbau, Segmentierung, Trigger] \href{https://youtu.be/NPGwFJqX0kc}{\textcolor{lpAccent}{https://youtu.be/NPGwFJqX0kc}}\\[2pt]
  {\footnotesize\color{lpGray}(URL: \nolinkurl{https://youtu.be/NPGwFJqX0kc})}
  \item[4.~Kampagnen und Automatisierung in der Praxis] \href{https://youtu.be/RF2m7mwJmxc}{\textcolor{lpAccent}{https://youtu.be/RF2m7mwJmxc}}\\[2pt]
  {\footnotesize\color{lpGray}(URL: \nolinkurl{https://youtu.be/RF2m7mwJmxc})}
\end{description}
\vspace{0.6em}
\noindent{\footnotesize\color{lpGray}\textit{Tipp:} \"Offnen Sie das Video, lassen Sie es im Hintergrund laufen und arbeiten Sie parallel an der Aufgabe.}

\clearpage

\aufgabenkopf{1}{Newsletter \textendash{} Kampagne \textendash{} Nurturing \textendash{} Vertriebsstrecke}{\nivLi}{8\,min}

\noindent Definieren Sie in jeweils \emph{einem} Satz die vier Begriffe, die im professionellen E-Mail-Vertrieb häufig vermischt werden.

\ytvideo{E-Mail-Marketing als systematischer Vertriebshebel}{https://youtu.be/GlKFLDx0Fjw}

\begin{loesungbox}
\begin{enumerate}
  \item \textbf{Newsletter:} regelmäßig versandte E-Mail an eine breite Empfängerliste, primär zur Information und Markenpflege, ohne enge Verbindung zu einer einzelnen Vertriebs- oder Kaufphase.
  \item \textbf{Kampagne:} thematisch und zeitlich begrenzte E-Mail-Aktion zu einem konkreten Anlass oder Angebot, mit definierter Zielgruppe, Zielzahl und Erfolgsmessung.
  \item \textbf{Lead-Nurturing:} systematische, mehrstufige E-Mail-Kommunikation, die Interessenten entlang der Customer Journey gezielt zur Kaufbereitschaft entwickelt \textendash{} Inhalte und Timing folgen dem Reifegrad des Leads.
  \item \textbf{Automatisierte Vertriebsstrecke:} regelbasierte, durch Verhalten ausgelöste E-Mail-Folge, die Lead-Nurturing technisch operationalisiert (Trigger, Bedingungen, Scoring, Übergabe an den Vertrieb).
  \item \textbf{Unterscheidung Newsletter vs.\ Nurturing-Strecke:} Der Newsletter ist \emph{kanalzentriert} (an alle, gleicher Zeitpunkt, gleicher Inhalt), die Nurturing-Strecke ist \emph{leadzentriert} (richtige Person, richtige Phase, richtiger Inhalt). Aus Vertriebssicht ist nur die Nurturing-Strecke ein steuerbarer Pipeline-Hebel.
\end{enumerate}
\end{loesungbox}

\clearpage

% --- AUFGABE 2 -----------------------------------------------------------
\aufgabenkopf{2}{E-Mail-Kennzahlen zuordnen}{\nivLi}{8\,min}

\noindent Ordnen Sie jede Kennzahl einer der drei Kategorien zu: (A) Aktivitäts-, (B) Qualitäts-, (C) wirtschaftliche Steuerungskennzahl.

\ytvideo{E-Mail-Marketing als systematischer Vertriebshebel}{https://youtu.be/GlKFLDx0Fjw}

\begin{loesungbox}
\begin{enumerate}
  \item \textbf{Öffnungsrate \textrightarrow{} (A) Aktivitätskennzahl.} Misst, ob die Betreffzeile zog \textendash{} kein Beleg für Vertriebswirkung.
  \item \textbf{Klickrate \textrightarrow{} (A) Aktivitätskennzahl.} Misst Interaktion mit dem Inhalt, aber nicht Kaufbereitschaft.
  \item \textbf{Conversion Rate zu Beratungstermin \textrightarrow{} (C) Wirtschaftliche Steuerungskennzahl.} Direkter Indikator dafür, dass die Strecke Vertriebsanlässe erzeugt.
  \item \textbf{Abmelderate \textrightarrow{} (B) Qualitätskennzahl.} Signal für Relevanz und Frequenz \textendash{} dauerhaft hohe Abmelderate beschädigt die Datenbasis.
  \item \textbf{Anteil vertriebsqualifizierter Leads (SQL) \textrightarrow{} (B) Qualitätskennzahl} mit wirtschaftlicher Konsequenz. Trennt ``Aktivität'' von ``Pipeline''.
  \item \textbf{Kosten pro qualifiziertem Lead \textrightarrow{} (C) Wirtschaftliche Steuerungskennzahl.} Erste echte ROI-Annäherung \textendash{} kombiniert Investition und Qualität.
  \item \textbf{Deckungsbeitrag aus E-Mail-Kampagne \textrightarrow{} (C) Wirtschaftliche Steuerungskennzahl.} Die finale Managementgröße, weil sie Marketing-Wirkung in Ergebnis übersetzt.
  \item \textbf{Bounce Rate \textrightarrow{} (B) Qualitätskennzahl.} Indikator für Datenhygiene, technische Hygiene und Listenpflege.
\end{enumerate}

\smallskip
\emph{Hinweis:} Öffnungs- und Klickrate sind nicht ``schlecht'' \textendash{} sie sind operative Diagnosegrößen. Steuerung auf Senior-Level läuft jedoch über Conversion, SQL-Anteil und ROI.
\end{loesungbox}

\clearpage

% --- AUFGABE 3 -----------------------------------------------------------
\aufgabenkopf{3}{ROI, CAC und Conversion Rate auf den Punkt}{\nivLi}{8\,min}

\noindent Erklären Sie ROI, CAC und Conversion Rate \textendash{} und ordnen Sie sie als Management- oder operative Steuerungskennzahl ein.

\ytvideo{E-Mail-Marketing als systematischer Vertriebshebel}{https://youtu.be/GlKFLDx0Fjw}

\begin{loesungbox}
\begin{enumerate}
  \item \textbf{ROI (Return on Investment):} setzt den wirtschaftlichen Ertrag einer Maßnahme (Deckungsbeitrag, Umsatz) ins Verhältnis zur eingesetzten Investition \textendash{} sagt: ``Lohnt sich diese Aktivität insgesamt?''
  \item \textbf{CAC (Customer Acquisition Cost):} durchschnittliche Kosten, einen \emph{neuen Kunden} zu gewinnen \textendash{} Summe aller Marketing- und Vertriebsausgaben geteilt durch Anzahl neu gewonnener Kunden.
  \item \textbf{Conversion Rate (zu Beratungstermin):} Anteil der Kontakte, die einen definierten nächsten Schritt tun \textendash{} sagt nur etwas über die Schritt-Effizienz, nicht über Wirtschaftlichkeit. Im Gegensatz dazu misst die Klickrate nur die \emph{Interaktion}, die Conversion Rate misst den \emph{Funnel-Übergang}.
  \item \textbf{Management vs.\ operativ:} \emph{ROI} ist die Management-Kennzahl (Geschäftsführung, CFO, CRO). \emph{CAC} steht zwischen operativ und strategisch \textendash{} er ist Steuerungsgröße für das Kanal-Portfolio. \emph{Conversion Rate} ist operativ \textendash{} sie steuert einzelne Kampagnen, Landingpages, E-Mail-Varianten.
\end{enumerate}
\end{loesungbox}

\clearpage

% =========================================================================
% L2 AUFGABEN
% =========================================================================
\section*{Niveau L2 \textendash{} Anwendung \& Transfer}
\addcontentsline{toc}{section}{L2 -- Anwendung}

% --- AUFGABE 4 -----------------------------------------------------------
\aufgabenkopf{4}{CloudDesk Solutions \textendash{} Nurturing-Strecke ohne Vorwissen}{\nivLii}{25\,min}

\begin{infobox}[Fallbeschreibung]
CloudDesk Solutions (Dokumentenmanagement-Software, Mittelstand): Whitepaper-Download wird oft gemacht, Beratungstermine folgen selten. Bisher: nur automatische Danke-Mail.
\end{infobox}

\ytvideo{Lead Nurturing mit Marketing Automation}{https://youtu.be/sjAkl8pwLgg}

\begin{loesungbox}
\textbf{1. Problem im aktuellen Prozess:}
Der Download liefert einen Lead-Datensatz \textendash{} aber zwischen Download und Beratungsanfrage liegt eine Lücke von Tagen bis Wochen, in der der Lead unbearbeitet bleibt. Die einzige Berührung ist eine Danke-Mail \textendash{} weder Vertrauensaufbau noch Argument für den nächsten Schritt. Das Resultat: hoher Aktivitätswert oben (Downloads), niedrige Conversion unten (Termine).

\smallskip
\textbf{2. Drei mögliche Gründe, warum Leads noch nicht buchen:}
\begin{itemize}
  \item \emph{Reifegrad:} Der Lead war auf der Suche nach Argumenten für seinen Chef, nicht nach einer Kaufentscheidung.
  \item \emph{Fehlendes Vertrauen:} Eine Datei reicht nicht, um den Anbieter als Lösungspartner einzuordnen \textendash{} kein Case, keine Referenz, keine Konkretisierung.
  \item \emph{Kein Anreiz für den Termin:} Der CTA ``Beratungstermin'' wirkt vertrieblich, nicht hilfreich \textendash{} es fehlt der Mehrwert für die ersten 15 Minuten.
\end{itemize}

\smallskip
\textbf{3.+4. E-Mail-Strecke (vier E-Mails) mit Ziel, Inhalt, CTA:}

\smallskip
\begin{tabularx}{\linewidth}{@{}lXXX@{}}
\toprule
\textbf{E-Mail} & \textbf{Ziel} & \textbf{Inhalt} & \textbf{CTA} \\
\midrule
1 (Tag 0) & Erwartungsmanagement, Vertrauen & Danke + 3 Highlights aus dem Whitepaper + ``Was kommt als Nächstes?'' & Whitepaper sicher ablegen / weiterleiten \\
\midrule
2 (Tag 3) & Problembewusstsein vertiefen & Kurze Case-Story: typische Dokumenten-Engpässe im Mittelstand & ``Lesen Sie die Praxisgeschichte'' (Blog) \\
\midrule
3 (Tag 7) & Lösung konkretisieren, Vertrauen festigen & Kunden-Referenz: Vorher / Nachher Zahlen einer ähnlichen Firma & ``Demo-Video ansehen (4 Min.)'' \\
\midrule
4 (Tag 14) & Vertriebskontakt anbieten & ``15 Minuten Sparringsgespräch'' \textendash{} nicht ``Verkaufstermin'' & ``Termin direkt buchen'' \\
\bottomrule
\end{tabularx}

\smallskip
\textbf{5. Übergabezeitpunkt an den Vertrieb:}
Erst wenn (a) Rolle erkennbar passt (Geschäftsführung, IT-, kaufmännische Leitung), (b) mindestens zwei E-Mail-Interaktionen über Klick/Antwort vorliegen \emph{oder} ein expliziter Termin-Klick erfolgt ist, (c) Unternehmensgröße im Zielsegment liegt. \emph{Faustregel:} Vertrieb übernimmt erst beim Demo-Klick, Video-Ansehen über 60\,\% oder direkter Anfrage \textendash{} nicht beim Whitepaper-Download.
\end{loesungbox}

\clearpage

% --- AUFGABE 5 -----------------------------------------------------------
\aufgabenkopf{5}{E-Mail-Performance analysieren und optimieren}{\nivLii}{25\,min}

\begin{infobox}[Kennzahlen]
1.200 Downloads, 62\,\% Öffnung Danke-Mail, 4\,\% Klick auf Beratung, 18 Termine (davon 7 qualifiziert), 3,8\,\% Abmeldung. Ziel: 40 qualifizierte Termine in 3 Monaten, Budget 18.000\,\euro, Vertriebskapazität max.\ 15 Termine/Woche.
\end{infobox}

\ytvideo{Lead Nurturing mit Marketing Automation}{https://youtu.be/sjAkl8pwLgg}

\begin{loesungbox}
\textbf{1. Bewertung der aktuellen Kampagne:}
\begin{itemize}
  \item \emph{Gut:} Lead-Generierung funktioniert (1.200 Downloads), Öffnungsrate 62\,\% ist hoch \textendash{} Interesse ist da.
  \item \emph{Mittelmäßig:} 4\,\% Klick auf Beratungslink ist niedrig, aber im B2B-Erstkontakt nicht ungewöhnlich.
  \item \emph{Kritisch:} Nur 7 von 18 Terminen sind qualifiziert \textendash{} also nur 39\,\% \emph{wirtschaftlich verwertbar}. 11 ``Termine'' kosten Vertriebszeit ohne Gegenwert.
  \item \emph{Verbesserungswürdig:} Abmelderate 3,8\,\% ist hoch \textendash{} Inhalt oder Frequenz ist nicht relevant genug.
\end{itemize}

\smallskip
\textbf{2. Vier Schwachstellen im Nurturing-Prozess:}
\begin{itemize}
  \item Es gibt nur \emph{eine} Mail \textendash{} kein Nurturing, sondern ein Pulsschlag.
  \item Keine Segmentierung nach Rolle, Unternehmensgröße oder Interesse \textendash{} ein Geschäftsführer und ein IT-Leiter brauchen unterschiedliche Inhalte.
  \item Übergabe an den Vertrieb erfolgt zu früh \textendash{} jeder Termin-Klicker landet im Kalender, unabhängig vom Reifegrad.
  \item Keine Re-Engagement-Logik für die 96\,\%, die nicht klicken \textendash{} sie verschwinden aus dem Funnel.
\end{itemize}

\smallskip
\textbf{3. Drei Optimierungsmaßnahmen:}
\begin{itemize}
  \item \emph{Inhalte:} 4-stufige Nurturing-Strecke (siehe Aufgabe 4) statt einer Danke-Mail.
  \item \emph{Segmentierung:} eine Variante für ``Geschäftsführung'' (ROI-Argumente), eine für ``IT-/kaufmännische Leitung'' (Integration, Datenschutz).
  \item \emph{Übergabe-Logik:} Lead Scoring (Rolle + Interaktion) \textendash{} Termin-Buchung erfordert mindestens 2 Klicks oder einen Demo-Video-Aufruf, sonst läuft der Lead weiter im Nurturing.
\end{itemize}

\smallskip
\textbf{4. Fünf KPIs für die neue Strecke:}
\begin{itemize}
  \item Conversion Rate zur Demo-Anfrage (operativ).
  \item Anteil vertriebsqualifizierter Termine (Qualität).
  \item Kosten pro qualifiziertem Termin (Wirtschaftlichkeit).
  \item Reifegrad zum Termin-Zeitpunkt (Score 0\,--\,100).
  \item Abmelderate je Variante (Listengesundheit).
\end{itemize}

\smallskip
\textbf{5. Ist das Ziel realistisch?}
40 qualifizierte Termine in 3 Monaten = ca.\ 13/Monat. Aktuell: 7 in einem Zyklus. Bei verdreifachter Datenbasis (bessere Re-Engagement-Strecke + zusätzliche Whitepaper-Downloads in 3 Monaten) und einer Verdopplung der \emph{Qualität} (Scoring statt jeder Klick) ist 40 \emph{ambitioniert, aber erreichbar} \textendash{} unter der Bedingung, dass das Vertriebsteam tatsächlich die strikte Übergabelogik akzeptiert. \emph{Risiko:} Wenn das Team nur Quantität sieht, wird Scoring umgangen und die Qualität bleibt bei 39\,\%. Empfehlung: Pilotphase 6 Wochen mit messbarem MQL/SQL-Check, dann Skalierung.
\end{loesungbox}

\clearpage

% --- AUFGABE 6 -----------------------------------------------------------
\aufgabenkopf{6}{Segmentierung einer Kontaktbasis}{\nivLii}{20\,min}

\noindent Entwickeln Sie eine Segmentierungslogik in vier Dimensionen.

\ytvideo{Lead Nurturing mit Marketing Automation}{https://youtu.be/sjAkl8pwLgg}

\begin{loesungbox}
\emph{Beispielhafte Segmentierungslogik \textendash{} andere Dimensionen sind legitim, solange sie operativ trennbar sind.}

\smallskip
\textbf{Dimension 1 \textendash{} Rolle / Buying-Center:}
\begin{itemize}
  \item Segmente: Geschäftsführung \textperiodcentered{} IT-Leitung \textperiodcentered{} Fachbereichsleitung \textperiodcentered{} Einkauf.
  \item Datenbedarf: Job-Titel-Feld im CRM, ggf.\ LinkedIn-Anreicherung.
  \item Inhalte: GF \textrightarrow{} ROI/Strategie \textperiodcentered{} IT \textrightarrow{} Integration/Sicherheit \textperiodcentered{} Fachbereich \textrightarrow{} Anwendungsfall \textperiodcentered{} Einkauf \textrightarrow{} Vertragsmodelle/TCO.
\end{itemize}

\textbf{Dimension 2 \textendash{} Branche:}
\begin{itemize}
  \item Segmente: Logistik \textperiodcentered{} produzierendes Gewerbe \textperiodcentered{} Dienstleister \textperiodcentered{} öffentliche Hand \textperiodcentered{} Gesundheit.
  \item Datenbedarf: Industriecode (NACE) oder manuelle Branchenkennzeichnung.
  \item Inhalte: branchenspezifische Cases (Logistik-Case wirkt nicht in Gesundheit).
\end{itemize}

\textbf{Dimension 3 \textendash{} Interesse / Themenpräferenz:}
\begin{itemize}
  \item Segmente: Kostenreduktion \textperiodcentered{} Compliance/Sicherheit \textperiodcentered{} Wachstum/Skalierung \textperiodcentered{} Prozessautomatisierung.
  \item Datenbedarf: Klick-Tracking auf Themenfelder, Whitepaper-Wahl, Webinar-Themen.
  \item Inhalte: Sprache und Argumentationslinie auf das jeweilige ``Motiv'' zuschneiden.
\end{itemize}

\textbf{Dimension 4 \textendash{} Reifegrad / Kaufphase:}
\begin{itemize}
  \item Segmente: Informationsphase \textperiodcentered{} Vergleichsphase \textperiodcentered{} Entscheidungsphase \textperiodcentered{} Bestandskunde.
  \item Datenbedarf: Lead-Score (Verhalten + Profil), Vertriebsfeedback.
  \item Inhalte: Information \textrightarrow{} Fachimpuls \textperiodcentered{} Vergleich \textrightarrow{} Checkliste/Wettbewerbsmatrix \textperiodcentered{} Entscheidung \textrightarrow{} Demo/Referenz \textperiodcentered{} Bestand \textrightarrow{} Erweiterung/Upgrade.
\end{itemize}

\smallskip
\textbf{Hypothese mit höchstem Hebel in 90 Tagen:} Das Segment ``IT- oder Fachbereichsleitung in der Vergleichsphase'' \textendash{} dort entscheidet sich die Shortlist, und genau hier liegen typischerweise die meisten Whitepaper- und Webinar-Leads ungenutzt im CRM. Eine gezielte Nurturing-Strecke ``Vergleichshilfe + Demo'' mit branchenpassendem Case erzeugt 90-Tage-Pipeline, ohne neue Lead-Quellen aufbauen zu müssen.
\end{loesungbox}

\clearpage

% --- AUFGABE 7 -----------------------------------------------------------
\aufgabenkopf{7}{EcoFleet Systems \textendash{} Nurturing- und KPI-Architektur}{\nivLii}{30\,min}

\begin{infobox}[Fallstudie: EcoFleet Systems]
Fuhrparkoptimierungs-Software. 8.000 Kontakte, 2.500 davon Webinar-/Whitepaper-Leads. Budget 95.000\,\euro{} für 6\,Monate. Ziel: 300 QL, 90 Demos, 25 Angebote, 8 Abschlüsse.
\end{infobox}

\ytvideo{Newsletter erfolgreich starten — Aufbau, Segmentierung, Trigger}{https://youtu.be/NPGwFJqX0kc}

\begin{loesungbox}
\textbf{1. Drei wichtigste Schwächen:}
\begin{itemize}
  \item Generischer Versand an alle 8.000 Kontakte \textendash{} keine Segmentierung nach Rolle, Branche oder Interesse, dadurch Relevanzverlust und steigende Abmelderate.
  \item Vertrieb erhält Leads zu früh \textendash{} keine MQL/SQL-Trennung, also Beladung mit unreifen Leads und Verzögerung bei den wirklich kaufbereiten.
  \item Erfolgsmessung über Öffnungs-/Klickraten \textendash{} keine Verbindung zwischen E-Mail-Aktivität und Pipeline-Beitrag, also keine ROI-Aussage möglich.
\end{itemize}

\smallskip
\textbf{2. Nurturing-Strecke (sechs Kontaktpunkte):}
\begin{itemize}
  \item \emph{E-Mail 1 (Tag 0):} Danke + Nutzenzusammenfassung des Whitepaper/Webinars.
  \item \emph{E-Mail 2 (Tag 4):} Problemimpuls \textendash{} typische Kostentreiber im Fuhrpark (rollenspezifisch).
  \item \emph{E-Mail 3 (Tag 10):} Praxisbeispiel \textendash{} Einsparpotenzial eines vergleichbaren Mittelstandsunternehmens.
  \item \emph{E-Mail 4 (Tag 17):} Fachinhalt \textendash{} CO\textsubscript{2}-Reporting und regulatorische Anforderungen (für GF/Nachhaltigkeit).
  \item \emph{E-Mail 5 (Tag 25):} Vergleichshilfe \textendash{} Checkliste für die Auswahl einer Fuhrparksoftware.
  \item \emph{E-Mail 6 (Tag 35):} Einladung zur Demo oder Erstberatung mit konkretem Nutzenversprechen (``In 30 Min.\ Ihr Einsparpotenzial berechnen'').
  \item \emph{E-Mail 7 optional (Tag 60):} Reaktivierung bei Inaktivität \textendash{} kurzer Entscheidungsleitfaden.
\end{itemize}

\smallskip
\textbf{3. Lead-Qualifizierung:}
\begin{itemize}
  \item \emph{MQL (Marketing Qualified Lead):} Zielgruppen-Rolle erfüllt + mindestens ein fachliches Interesse erkennbar + Einwilligung vorhanden + mindestens eine Interaktion über die Danke-Mail hinaus.
  \item \emph{SQL (Sales Qualified Lead):} konkrete Rolle im Buying-Center + erkennbarer Bedarf (z.\,B.\ Demo-Klick, Vergleich-Checkliste vollständig durchgegangen) + passende Unternehmensgröße + wiederholte Interaktion (Score $\geq 60$).
  \item Vertrieb übernimmt \emph{erst} SQL \textendash{} MQL bleibt im Marketing.
\end{itemize}

\smallskip
\textbf{4. KPI-Architektur (sechs Kennzahlen):}
\begin{itemize}
  \item Zustellrate (Datenqualität).
  \item Klickrate je Segment (Relevanz).
  \item Conversion Rate zu Demo (Funnel-Übergang).
  \item Anteil SQL an MQL (Qualifizierungseffizienz).
  \item Demo-zu-Angebot-Quote (Vertriebswirkung).
  \item Deckungsbeitrag je Kampagnenstrecke (ROI).
\end{itemize}
\end{loesungbox}

\clearpage

% --- AUFGABE 8 -----------------------------------------------------------
\aufgabenkopf{8}{Budgetverteilung 95.000\,\euro{} \textendash{} mit Begründung}{\nivLii}{20\,min}

\ytvideo{Newsletter erfolgreich starten — Aufbau, Segmentierung, Trigger}{https://youtu.be/NPGwFJqX0kc}

\begin{loesungbox}
\textbf{Beispielhafte Budgetverteilung:}

\smallskip
\begin{tabularx}{\linewidth}{@{}lXl@{}}
\toprule
\textbf{Posten} & \textbf{Begründung} & \textbf{Betrag} \\
\midrule
Contententwicklung & Cases, Checklisten, Fachimpulse sind das ``Werkzeug'' jeder Nurturing-Strecke \textendash{} ohne Inhalt keine Conversion. & 25.000\,\euro \\
\midrule
Marketing-Automation und technische Einrichtung & Tool-Setup, Trigger-Logik, Templates \textendash{} ohne saubere Mechanik keine Skalierung. & 20.000\,\euro \\
\midrule
CRM-Datenbereinigung und Segmentierung & Ohne saubere Daten ist Segmentierung wertlos und Scoring falsch \textendash{} stille, aber kritische Investition. & 15.000\,\euro \\
\midrule
Landing Pages, Demo-Formulare, Conversion-Optimierung & Schnittstelle vom Klick zum Termin \textendash{} hier entscheidet sich die wirtschaftliche Wirkung. & 15.000\,\euro \\
\midrule
A/B-Tests & Geringes Stückbudget, hoher Lerneffekt \textendash{} senkt Risiko falscher Optimierungsentscheidungen. & 10.000\,\euro \\
\midrule
Reporting und ROI-Auswertung & Macht das System Senior-Level-fähig \textendash{} ohne sichtbares ROI keine Budgetabsicherung im nächsten Jahr. & 10.000\,\euro \\
\midrule
\textbf{Summe} & & \textbf{95.000\,\euro} \\
\bottomrule
\end{tabularx}

\smallskip
\textbf{Empfehlung in drei Sätzen:}
Die zwei wichtigsten Posten sind \emph{Contententwicklung} und \emph{CRM-Datenbereinigung} \textendash{} der eine entscheidet über Relevanz, der andere über Steuerbarkeit, beide gemeinsam über jeden weiteren Hebel. Reduziert werden würde am ehesten der A/B-Test-Posten \textendash{} weil A/B-Tests ohne ausreichendes Volumen ohnehin nur wenige Iterationen erlauben. Marketing-Automation und Reporting sind nicht-verhandelbar, weil sie das System überhaupt erst messbar machen.
\end{loesungbox}

\clearpage

% =========================================================================
% L3 AUFGABEN
% =========================================================================
\section*{Niveau L3 \textendash{} Management \& Entscheidungsarchitektur}
\addcontentsline{toc}{section}{L3 -- Management}

% --- AUFGABE 9 -----------------------------------------------------------
\aufgabenkopf{9}{EcoFleet \textendash{} Entscheidung unter Unsicherheit}{\nivLiii}{40\,min}

\begin{infobox}[Rolle und Spannungsfeld]
CRO von EcoFleet. Marketing will mehr Leads, Vertrieb will bessere Leads, Geschäftsführung will Zahlen. Entscheidung: Quantität zuerst oder Qualität zuerst?
\end{infobox}

\ytvideo{Newsletter erfolgreich starten — Aufbau, Segmentierung, Trigger}{https://youtu.be/NPGwFJqX0kc}

\begin{loesungbox}
\textbf{1. Analyse des Spannungsfelds:}
\begin{itemize}
  \item \emph{Marketing} denkt in Reichweite, weil Wachstum sichtbarer ist als Effizienz. Annahme: ``Mehr Volumen wäscht den Trichter selbst.''
  \item \emph{Vertrieb} denkt in Kapazität, weil unreife Leads echte Opportunitätskosten erzeugen. Annahme: ``Bestandsleads sind besser als ihr Ruf, wenn sie richtig bedient würden.''
  \item \emph{Geschäftsführung} denkt in Abschlüssen, weil sie das Budget gegen einen wirtschaftlichen Beitrag rechtfertigen muss.
  \item Hinter dem Konflikt steht ein Konstruktionsfehler: Es gibt keine gemeinsame Definition von ``Lead'' und keine gemeinsame Sicht auf den Pipeline-Beitrag.
  \item Beide Seiten haben recht \textendash{} aber nur eine Reihenfolge ist wirtschaftlich.
\end{itemize}

\smallskip
\textbf{2. Zwei strategische Optionen:}
\begin{itemize}
  \item \emph{Option A \textendash{} Quantität zuerst:} 60\,k\,\euro{} in neue Lead-Quellen (LinkedIn-Ads, zusätzliches Webinar, Whitepaper-Kampagne), 35\,k\,\euro{} in Tooling/Reporting. Bestandsdatenbank bleibt unverändert.
  \item \emph{Option B \textendash{} Qualität zuerst:} 55\,k\,\euro{} in Segmentierung, CRM-Bereinigung, Nurturing-Inhalte, MQL/SQL-Definition für die 2.500 bestehenden Leads. 40\,k\,\euro{} in Tooling, Reporting, Landing Pages. Neue Lead-Quellen erst ab Monat 4 selektiv.
\end{itemize}

\smallskip
\textbf{3. Bewertung beider Optionen:}

\smallskip
\begin{tabularx}{\linewidth}{@{}lXX@{}}
\toprule
\textbf{Kriterium} & \textbf{A: Quantität} & \textbf{B: Qualität} \\
\midrule
Zielerreichung 300\,QL & wahrscheinlich (Volumen) & wahrscheinlich (Conversion-Hub) \\
Risiko Vertriebskapazität & hoch \textendash{} mehr unreife Leads & niedrig \textendash{} weniger, aber bessere Leads \\
ROI-Sichtbarkeit nach 6\,M.\ & mittel \textendash{} Aktivitäten messbar, Qualität unklar & hoch \textendash{} Pipeline-Beitrag je Segment isolierbar \\
Datenkontrolle & gleich & deutlich besser \\
Skalierbarkeit ab Monat 7 & ungerichtet \textendash{} mehr vom Gleichen & gerichtet \textendash{} man weiß, welcher Kanal welche Conversion liefert \\
\bottomrule
\end{tabularx}

\smallskip
\textbf{4. Entscheidung und Budget:}
\emph{Qualität zuerst, dann selektive Skalierung.}
\begin{itemize}
  \item 55\,k\,\euro{} Qualität: Segmentierung (15), Content (20), CRM-Hygiene (10), MQL/SQL-Definition + Schulung (10).
  \item 25\,k\,\euro{} Tooling, Landing Pages, A/B-Tests.
  \item 15\,k\,\euro{} Reporting und Pilot-Auswertung.
  \item Erst ab Monat 4: bis zu 30\,k\,\euro{} (aus dem Folgebudget) gezielt in den \emph{einen} Kanal, der nach Pilot die beste SQL-Quote liefert.
\end{itemize}

\smallskip
\textbf{5. Pilot-Logik (erste 8 Wochen):}
\begin{itemize}
  \item Zwei Pilot-Segmente: ``Fuhrparkleitung Logistik'' und ``GF Mittelstand Bau''.
  \item Je 250 Leads aus dem Bestand bekommen die neue 6-stufige Nurturing-Strecke.
  \item Steuergröße: SQL-Anteil und Demo-zu-Angebot-Quote.
  \item Korrekturpunkt: Wenn nach 8 Wochen der SQL-Anteil nicht mindestens 12\,\% (gegen aktuell 39\,\% Qualität \emph{nach} Termin, was hier auf MQL-Stufe übersetzt $\approx$ 10\,--\,12\,\%) erreicht, wird die Inhalts-Logik nachjustiert \emph{vor} der Skalierung.
\end{itemize}

\smallskip
\textbf{6. Entscheidung unter Unsicherheit:}
Die Entscheidung ``Qualität zuerst'' wird getroffen, obwohl unklar ist, ob die 2.500 Bestandsleads tatsächlich noch entwickelbar sind \textendash{} oder ob sie inhaltlich bereits ``durchgespielt'' wurden. \emph{Begründung trotz Unsicherheit:}
\begin{itemize}
  \item Die Kosten dieser Entscheidung sind sichtbar (Investition in Inhalte, Hygiene), die Kosten der Gegenentscheidung (überlasteter Vertrieb, sinkende Abschlussquote, demotiviertes Team) wären verzögert sichtbar \textendash{} und teurer.
  \item Wenn die Annahme falsch ist, lässt sich Qualität in 6 Wochen messen und korrigieren; eine falsche Quantitätsentscheidung lässt sich erst nach Monaten an der Abschlussquote ablesen.
  \item Bei knapper Vertriebskapazität ist die Reihenfolge \emph{erst Qualität, dann Quantität} die einzige, die nicht durch Folgekosten gefährdet ist.
\end{itemize}
\end{loesungbox}

\clearpage

% --- AUFGABE 10 ----------------------------------------------------------
\aufgabenkopf{10}{Zielkonflikt: Automatisierung vs.\ Personalisierung}{\nivLiii}{30\,min}

\noindent\textbf{Diskussionsaufgabe.}

\ytvideo{Kampagnen und Automatisierung in der Praxis}{https://youtu.be/RF2m7mwJmxc}

\begin{loesungbox}
\textbf{1. Wann kippt Automatisierung in Belästigung:}
Automatisierung kippt typischerweise an drei Schwellen. \emph{Erstens}, wenn Frequenz ohne sichtbaren Nutzen steigt \textendash{} der Empfänger lernt, dass die Mails ``Rauschen'' sind. \emph{Zweitens}, wenn personalisierte Felder (Vorname, Firma) falsch oder offensichtlich automatisiert wirken \textendash{} Pseudo-Persönlichkeit beschädigt Vertrauen stärker als ein offen massenhafter Newsletter. \emph{Drittens}, wenn Verhalten zwar getrackt, aber nicht respektvoll verarbeitet wird \textendash{} etwa wenn ein Lead Verträge unterzeichnet hat und trotzdem weiter mit ``Erstkontakt''-Inhalten bespielt wird. Kritische Schwelle ist außerdem: Mehr als zwei E-Mails pro Woche ohne klaren Nutzen führen messbar zu Abmelderaten über 4\,\%.

\smallskip
\textbf{2. Konkretes B2B-Szenario:} SaaS-Anbieter mit 50.000 Kontakten und Bestandskunden.

\smallskip
\emph{Drei sinnvolle Automatisierungen:}
\begin{itemize}
  \item Trigger-basierte Welcome-Strecke nach Trial-Anmeldung \textendash{} hoher Mehrwert, klare Kundenerwartung.
  \item Inaktivitäts-Reaktivierung (nach 90 Tagen ohne Login) mit konkretem Hilfsangebot \textendash{} der Kunde verliert sonst sowieso den Vertrag.
  \item Vertragsverlängerungs-Strecke 60/30/14 Tage vor Ablauf \textendash{} Service-Charakter, nicht Vertriebsdruck.
\end{itemize}

\emph{Eine Automatisierung, die eine Grenze überschreitet:}
\begin{itemize}
  \item ``Automatisch jeden Lead, der dreimal die Pricing-Seite besucht hat, an einen Vertriebler übergeben mit Hinweis: \emph{Hat Kaufabsicht, sofort anrufen}.'' \textendash{} verletzt Datensensibilität, erzeugt Anrufdruck und beschädigt Vertrauen.
\end{itemize}

\emph{Drei realistische Personalisierungs-Hebel:}
\begin{itemize}
  \item Rollenbasierte Variante einer Mail (GF / IT / Fachbereich) statt 47 Mikro-Personalisierungen.
  \item Branchenbasierter Case (ein einziger guter Case je Branche, sauber gepflegt) statt zehn halbgaren.
  \item Reifegrad-basierter Tonfall (Information $\rightarrow$ Vergleich $\rightarrow$ Entscheidung) statt jeder Mail ein Verkaufsabschluss.
\end{itemize}

\emph{Konflikte:}
\begin{itemize}
  \item \emph{Datenschutz:} Verhaltensbasierte Personalisierung erfordert ausdrückliche Einwilligung; ohne sie ist Trigger-Logik rechtlich nicht haltbar.
  \item \emph{Vertriebskapazität:} Personalisierung erzeugt mehr Konversation, die das Team binden kann \textendash{} ohne Kapazitätsplanung wird ``personalisiert'' zu ``unbeantwortet''.
  \item \emph{Markenwahrnehmung:} Hochpersonalisierte Mails wirken in einigen Branchen (Industrie, Gesundheit) übergriffig \textendash{} dort ist nüchtern überpersönlich.
\end{itemize}

\smallskip
\textbf{3. Faustregel:}
Automatisierung ist wirtschaftlich begründet, solange sie (a) einen für den Empfänger \emph{nachvollziehbaren} Anlass hat \emph{und} (b) die Personalisierung an realen Daten und nicht an Pseudo-Variablen hängt. Sie beginnt Vertrauen zu zerstören, sobald der Empfänger den Eindruck hat, als ``Datensatz'' und nicht als Person behandelt zu werden \textendash{} und das ist meist deutlich \emph{vor} dem Punkt, an dem die Abmelderate auffällt.
\end{loesungbox}

\clearpage

% --- AUFGABE 11 ----------------------------------------------------------
\aufgabenkopf{11}{Transferprojekt \textendash{} E-Mail-, Nurturing- \& ROI-Architektur}{\nivLiii}{45\,min}

\begin{infobox}[Rolle und Ausgangslage]
CRO / Head of Marketing Automation / Leiter:in Digital Sales eines mittelständischen B2B-Unternehmens mit wachsender Kontaktdatenbank, ohne Nurturing-Architektur, ohne ROI-Messung, mit Spannungen zwischen Marketing und Vertrieb.
\end{infobox}

\ytvideo{Kampagnen und Automatisierung in der Praxis}{https://youtu.be/RF2m7mwJmxc}

\begin{loesungbox}
\emph{Musterlösung als Entscheidungsarchitektur \textendash{} andere Konfigurationen sind möglich, solange sie als Architektur (nicht als Maßnahmenliste) gedacht sind.}

\smallskip
\textbf{1. Zielbild:}
Bis Ende des Folgejahres ist E-Mail-Marketing kein Versandkanal mehr, sondern ein \emph{Pipeline-Steuerungssystem}. Jede Strecke ist segmentiert, jeder Lead hat einen Score, jeder Übergabepunkt an den Vertrieb ist definiert. Die Geschäftsführung bekommt monatlich einen ROI-Bericht je Kampagnenstrecke. Marketing wird an Pipeline-Beitrag (nicht an Versandzahlen) gemessen.

\smallskip
\textbf{2. Segmentierung der Kontaktbasis:}
\begin{itemize}
  \item \emph{Nach Rolle:} GF \textperiodcentered{} IT \textperiodcentered{} Fachbereich \textperiodcentered{} Einkauf.
  \item \emph{Nach Branche:} drei bis fünf Kernbranchen, separate Inhalte je Branche.
  \item \emph{Nach Interesse:} Kosten, Compliance, Wachstum, Automatisierung.
  \item \emph{Nach Verhalten:} Whitepaper-Download, Webinar, Demo-Anfrage, Inaktivität.
  \item \emph{Nach Reifegrad:} Information \textperiodcentered{} Vergleich \textperiodcentered{} Entscheidung \textperiodcentered{} Bestand.
\end{itemize}

\smallskip
\textbf{3. Nurturing-Logik entlang der Customer Journey:}
Erstkontakt (Danke + Erwartung) \textrightarrow{} Problembewusstsein (Impuls) \textrightarrow{} Lösungsvergleich (Case + Checkliste) \textrightarrow{} Entscheidung (Demo / Beratung) \textrightarrow{} Übergabe an Vertrieb \textrightarrow{} Onboarding \textrightarrow{} Bestandskunden-Strecke (Upsell, Reaktivierung).

\smallskip
\textbf{4. Inhaltsarchitektur:}
\begin{itemize}
  \item E-Mail-Typen: Welcome, Fachimpuls, Case, Webinar-Einladung, Vergleichshilfe, Demo-Einladung, Reaktivierung, Vertragsverlängerung, Upsell.
  \item Trigger: Download, Webinar-Anmeldung, Websitebesuch ($\geq 3$ Mal), Inaktivität ($> 60$\,Tage), Demo-Klick, Angebotsablauf.
  \item CTAs: ``Lesen Sie den Case'', ``Buchen Sie 15 Min.\ Sparring'', ``Vergleich starten'', ``Demo ansehen''.
\end{itemize}

\smallskip
\textbf{5. Lead-Scoring und Übergabe:}
\begin{itemize}
  \item Profilscore (max.\ 40): Rolle, Unternehmensgröße, Branche.
  \item Verhaltensscore (max.\ 60): Klicks, Demo-Klick, Downloads, Webinar-Teilnahme.
  \item MQL ab Gesamtscore 40, SQL ab 70 \emph{oder} expliziter Termin-Klick mit passendem Profil.
  \item Übergabe SQL an Vertrieb innerhalb von 24\,Stunden, mit kompaktem Lead-Briefing aus dem CRM.
\end{itemize}

\smallskip
\textbf{6. Daten- und Datenschutzlogik:}
\begin{itemize}
  \item Single Opt-in für allgemeine Inhalte, Double Opt-in mit Zweck-Trennung für Verhaltensbasiertes.
  \item CRM ist System of Truth \textendash{} alle Tools schreiben dorthin, nicht umgekehrt.
  \item Datenhygiene-Routine: jährliche Bereinigung inaktiver Kontakte (kein Versand $> 12$\,Monate \textrightarrow{} aktive Reaktivierung oder Löschung).
  \item DSGVO-Dokumentation: Einwilligungsbeleg, Speicherzweck, Speicherdauer, Auskunftsprozess sind formal definiert.
\end{itemize}

\smallskip
\textbf{7. KPI-Architektur:}
\begin{itemize}
  \item \emph{Operativ:} Zustellrate, Bounce Rate, Öffnungs-/Klickrate je Segment.
  \item \emph{Qualität:} MQL-Anteil, SQL-Anteil, Abmelderate, Reifegrad-Verteilung.
  \item \emph{Wirtschaft:} CAC, Kosten je SQL, Demo-zu-Angebot-Quote, Abschluss-Quote, Umsatz und Deckungsbeitrag je Strecke.
\end{itemize}

\smallskip
\textbf{8. ROI-Logik:}
\begin{itemize}
  \item ROI je Strecke = (Deckungsbeitrag aus SQL-generierten Abschlüssen $-$ Streckenkosten) / Streckenkosten.
  \item Bewertung je Segment, je Kanalquelle, je Kampagnenstrecke \textendash{} keine ROI-Bewertung ohne Verbindung zum tatsächlichen Vertriebsergebnis.
  \item Reporting-Rhythmus: monatlich operativ, quartalsweise auf Geschäftsleitungs-Ebene.
\end{itemize}

\smallskip
\textbf{9. Governance:}
\begin{itemize}
  \item \emph{CRO} entscheidet über Kanal-Portfolio, MQL/SQL-Definition und Pipeline-Ziele.
  \item \emph{Head of Marketing Automation} verantwortet operative Streckenplanung, Tool-Logik, Inhalte.
  \item \emph{Vertriebsleitung} verantwortet Lead-Annahme, Feedback-Loop, SLA für Lead-Bearbeitung.
  \item \emph{IT / Datenschutzbeauftragter} verantworten CRM, Einwilligung, Datenschutz-Compliance.
  \item \emph{Controlling} verantwortet ROI-Logik und Kennzahlen-Definition.
  \item Gemeinsames Forum: ``Revenue Council'' monatlich \textendash{} ein Beschluss je Quartal, keine Diskussionsrunde.
\end{itemize}

\smallskip
\textbf{10. Drei Managemententscheidungen unter Unsicherheit:}
\begin{enumerate}
  \item \emph{MQL/SQL-Trennung wird verbindlich eingeführt, bevor das Vertriebsteam das Konzept akzeptiert.} Begründung: ohne Definition kein Score, ohne Score keine Steuerung; Akzeptanz folgt der Wirkung, nicht umgekehrt.
  \item \emph{Erste 30\,\% des Budgets fließen in CRM-Datenqualität, nicht in neue Lead-Quellen.} Begründung: die Investition ist intern unsichtbar, aber jede weitere Maßnahme darauf angewiesen.
  \item \emph{Reporting wird quartalsweise auf Geschäftsleitungs-Ebene gezogen, bevor die Zahlen ``schön'' sind.} Begründung: transparente schwache Zahlen sind besser als verzögerte gute \textendash{} Vertrauen in die Marketingsteuerung entsteht durch Konsequenz, nicht durch Kosmetik.
\end{enumerate}

\smallskip
\textbf{Zusatz \textendash{} Entscheidung gegen mehr Kampagnenvolumen:}
Bewusst \emph{nicht} eingeführt wird ein zusätzlicher monatlicher Massen-Newsletter an die gesamte Datenbank von 50.000+ Kontakten. Aus Senior-Level-Perspektive: jeder zusätzliche Versand ohne Reifegrad-Bezug schadet der Listengesundheit (Abmelderate, Bounce, Spam-Reputation), bindet Content-Kapazität ohne nachweisbaren Pipeline-Beitrag und belastet das Vertriebsteam mit unreifen Reaktionen. Stattdessen: gleicher Content-Aufwand fließt in zwei rollenbasierte Nurturing-Strecken \textendash{} weniger ``Versand'', mehr Wirkung.
\end{loesungbox}

\clearpage

% --- Abschluss -----------------------------------------------------------
\section*{Abschluss}
\thispagestyle{plain}

\noindent Die Musterlösungen sind eine Diskussionsbasis. Wenn Ihre eigene Lösung anders ist, fragen Sie: \emph{Habe ich andere Annahmen \textendash{} oder einen anderen Maßstab angelegt?} Beides ist legitim, solange die Logik nachvollziehbar bleibt.

\medskip
\begin{infobox}[Nächster Schritt]
Bringen Sie Ihre eigenen Lösungen in die Coaching-Sitzung mit. Im Fokus stehen drei Fragen: Welche Annahmen haben Sie gemacht? Welche Zielkonflikte haben Sie sichtbar gemacht? Welche Entscheidung haben Sie getroffen \textendash{} und würden Sie auch verantworten, wenn sich der Markt anders entwickelt als angenommen?
\end{infobox}

\vspace{1cm}
\noindent{\color{lpAccent}\rule{\linewidth}{0.6pt}}\\[4pt]
{\footnotesize\sffamily\color{lpGray}Lernplatt \textperiodcentered{} by Can Siebert \textperiodcentered{} Kurs 22 (Bo 143) \textperiodcentered{} Tag 10 \textperiodcentered{} Lösungsheft \textperiodcentered{} \textcopyright{} 2026}

\end{document}
